\chapter[Charakterystyki i PDE pierwszego rzędu]{Równania pierwszego rzędu i metoda charakterystyk}
\label{chap:pde-characteristics}

W tym rozdziale zamienimy równania pierwszego rzędu na układy
równań zwyczajnych. Wersja liniowa metody pozwala rozwiązywać
równania transportu i ciągłości, natomiast wersja quasiliniowa
wyjaśnia powstawanie osobliwości z gładkich danych. Dla równań
w pełni nieliniowych skonstruujemy układ charakterystyczny
Charpita i zastosujemy go do równań Hamiltona--Jacobiego.

W całym rozdziale $d\geq1$, $T>0$ oraz $\lambda_d$ oznacza
miarę Lebesgue'a na $\R^d$. Wszystkie pochodne występujące
w sformułowaniach klasycznych rozumiemy w sensie punktowym.
W szczególności przez $C^1([0,T]\times\R^d)$ rozumiemy
funkcje o pochodnych pierwszego rzędu ciągłych aż do
$t=0,T$ (jednostronnie w tych punktach).

\section{Pochodna wzdłuż krzywej i najprostsze równanie transportu}
\label{sec:pde2-basic}

Niech $u\in C^1((0,T)\times\R^d)$ i niech
$X\colon I\to\R^d$ będzie krzywą klasy $C^1$, gdzie
$I\subset(0,T)$ jest przedziałem. Reguła łańcuchowa daje
\begin{equation}\label{eq:pde2-chain}
 \frac{\mathrm d}{\mathrm dt}u(t,X(t))
 =u_t(t,X(t))+\nabla_xu(t,X(t))\cdot\dot X(t).
\end{equation}
Dobierając $\dot X(t)$ do współczynników równania,
możemy przekształcić PDE w równanie zwyczajne
wzdłuż krzywej $t\mapsto(t,X(t))$. Taką krzywą nazywamy
\emph{charakterystyką} równania transportu.

\begin{stwierdzenie}[Transport ze stałą prędkością i źródłem]
\label{prop:pde2-constant}
Niech $c\in\R^d$, $u_0\in C^1(\R^d)$,
$f\in C^1([0,T]\times\R^d)$. Istnieje dokładnie jedno
rozwiązanie $u\in C^1([0,T]\times\R^d)$ zagadnienia
\[
 u_t+c\cdot\nabla_xu=f(t,x),\qquad u(0,x)=u_0(x).
\]
Jest ono dane wzorem
\begin{equation}\label{eq:pde2-constant-formula}
 u(t,x)=u_0(x-tc)+\int_0^t f(s,x-(t-s)c)\,\mathrm d\lambda_1(s).
\end{equation}
\end{stwierdzenie}
\begin{proof}
Charakterystyka wychodząca z $\xi$ ma postać
$X(t;\xi)=\xi+tc$. Dla każdego rozwiązania klasy $C^1$
reguła \eqref{eq:pde2-chain} pokazuje, że
$\frac{\mathrm d}{\mathrm dt}u(t,\xi+tc)=f(t,\xi+tc)$.
Całkując od $0$ do $t$ i podstawiając $\xi=x-tc$,
otrzymujemy \eqref{eq:pde2-constant-formula}; dowodzi to
jednoznaczności. Odwrotnie, oznaczmy prawą stronę wzoru
\eqref{eq:pde2-constant-formula} przez $u(t,x)$.
Na każdym zwartym otoczeniu parametrów $(t,x)$ funkcja
podcałkowa oraz jej pochodne względem $t$ i $x$ są ciągłe
i ograniczone dla $s\in[0,T]$. Reguła Leibniza dla całki
z parametrem daje więc ciągłe pochodne pierwszego rzędu,
także jednostronnie na końcach przedziału czasu.
Po podstawieniu $x=\xi+tc$ otrzymujemy
$u(t,\xi+tc)=u_0(\xi)+\int_0^t f(s,\xi+sc)\,\mathrm d\lambda_1(s)$.
Różniczkując względem $t$, dostajemy
$u_t(t,\xi+tc)+c\cdot\nabla_xu(t,\xi+tc)=f(t,\xi+tc)$.
Ponieważ dla ustalonego $t$ przesunięcie $\xi\mapsto\xi+tc$
jest bijekcją, równanie zachodzi dla każdego $x\in\R^d$.
Dla $t=0$ całka znika, więc $u(0,x)=u_0(x)$.
\end{proof}

\begin{przyklad}[Tłumienie wzdłuż charakterystyk]
Niech $\gamma>0$. Dla równania
$u_t+c\cdot\nabla u+\gamma u=0$ z danymi $u_0$
otrzymujemy $u(t,x)=e^{-\gamma t}u_0(x-tc)$.
Transport przesuwa profil, podczas gdy czynnik wykładniczy
tłumi jego amplitudę.
\end{przyklad}

\section{Zmienne pole prędkości: przepływ charakterystyk}
\label{sec:pde2-linear}

Rozważmy równanie
\begin{equation}\label{eq:pde2-linear}
 u_t+b(t,x)\cdot\nabla_xu+c(t,x)u=f(t,x)
 \quad\text{na }(0,T)\times\R^d.
\end{equation}
Dla przejrzystości przyjmiemy dość mocne założenia
zapewniające globalny przepływ; później będziemy je osłabiać.

\begin{twierdzenie}[Metoda charakterystyk dla transportu liniowego]
\label{thm:pde2-linear}
Niech $b\colon[0,T]\times\R^d\to\R^d$ będzie klasy $C^1$
i niech istnieje $M>0$ takie, że
$\|D_xb(t,x)\|\leq M$ i $|b(t,x)|\leq M(1+|x|)$
dla wszystkich $(t,x)\in[0,T]\times\R^d$, gdzie
$\|A\|=\sup_{|v|=1}|Av|$ oznacza normę operatorową
indukowaną przez normę euklidesową. Niech $c,f\in C^1([0,T]\times\R^d)$
oraz $u_0\in C^1(\R^d)$. Wówczas istnieje dokładnie jedno
rozwiązanie $u\in C^1([0,T]\times\R^d)$ zagadnienia
\eqref{eq:pde2-linear} z $u(0,x)=u_0(x)$.

Niech $X(t;\xi)$ będzie rozwiązaniem
\begin{equation}\label{eq:pde2-flow-ode}
 \dot X(t;\xi)=b(t,X(t;\xi)),\qquad X(0;\xi)=\xi.
\end{equation}
Dla każdego $t\in[0,T]$ odwzorowanie
$X_t\colon\xi\mapsto X(t;\xi)$ jest dyfeomorfizmem klasy $C^1$
przestrzeni $\R^d$ na siebie. Rozwiązanie jest określone przez
\begin{equation}\label{eq:pde2-solution}
 \begin{aligned}
 u(t,X(t;\xi))={}&e^{-C(t;\xi)}u_0(\xi)\\
 &+\int_0^t e^{-(C(t;\xi)-C(s;\xi))}
 f(s,X(s;\xi))\,\mathrm ds,
 \end{aligned}
\end{equation}
gdzie $C(t;\xi)=\int_0^t c(r,X(r;\xi))\,\mathrm dr$.
\end{twierdzenie}
\begin{proof}
Pole $b$ jest lokalnie lipschitzowskie względem $x$,
a dzięki warunkowi na $D_xb$ jest jednostajnie
lipschitzowskie na całej przestrzeni. Na przedziale istnienia rozwiązania jego równanie
całkowe i warunek wzrostu dają
\[
 1+|X(t;\xi)|\le 1+|\xi|
 +M\int_0^t(1+|X(s;\xi)|)\,\mathrm d\lambda_1(s).
\]
Z nierówności Gronwalla wynika
$1+|X(t;\xi)|\le(1+|\xi|)e^{Mt}$.
Gdyby maksymalny przedział istnienia kończył się przed $T$,
trajektoria pozostawałaby w ustalonej kuli domkniętej.
Pole $b$ jest ograniczone na iloczynie tej kuli i $[0,T]$,
więc $X$ jest lipschitzowska względem czasu i ma granicę
w końcu przedziału istnienia. Lokalne twierdzenie o istnieniu
ODE z tym warunkiem początkowym przedłuża trajektorię,
co przeczy maksymalności. To samo ograniczenie zapewnia
istnienie granicy przy $t=T$; równanie całkowe i ciągłość
pola dają jednostronną pochodną w tym punkcie. Twierdzenie o zależności rozwiązań
ODE od warunków początkowych daje $C^1$-gładkość $X$
względem $(t,\xi)$.

Dla ustalonego $t\in[0,T]$ i $x\in\R^d$ rozważmy
równanie $\dot Y(s)=-b(t-s,Y(s))$, $Y(0)=x$,
dla $s\in[0,t]$. Pole po prawej stronie spełnia te same
ograniczenia wzrostu i pochodnej przestrzennej co $b$,
więc poprzedni argument zapewnia istnienie na całym
$[0,t]$. Krzywa $r\mapsto Y(t-r)$ rozwiązuje pierwotne
równanie na $[0,t]$. Z jednoznaczności wynika
$X(t;Y(t))=x$ oraz, jeśli $x=X(t;\xi)$, to $Y(t)=\xi$.
Odwrotność $X_t^{-1}$ ma zatem postać $x\mapsto Y(t)$
i jest klasy $C^1$ z twierdzenia o zależności od danych.
W szczególności $D_\xi X(t;\xi)$ jest odwracalna.

Rozważmy odwzorowanie $\Phi(t,\xi)=(t,X(t;\xi))$.
Jego macierz pochodnej ma postać
\[
 D\Phi(t,\xi)=
 \begin{pmatrix}1&0\\ \partial_t X(t;\xi)&D_\xi X(t;\xi)\end{pmatrix},
\]
więc jej wyznacznik jest niezerowy.
Wewnątrz przedziału czasu twierdzenie o funkcji odwrotnej
zapewnia lokalną odwrotność klasy $C^1$. Ponieważ każde $X_t$
jest bijekcją, lokalne odwrotności są zgodne i dają wspólną
odwrotność $\Phi^{-1}(t,x)=(t,X_t^{-1}(x))$.
Na końcach przedziału stosujemy lokalnie to samo twierdzenie
do przedłużenia $X$ klasy $C^1$ względem czasu:
dla $t<0$ przyjmujemy $2X(0;\xi)-X(-t;\xi)$, a dla $t>T$
przyjmujemy $2X(T;\xi)-X(2T-t;\xi)$.
Pochodne czasowe i przestrzenne obu przedłużeń zgadzają się
z odpowiednimi jednostronnymi pochodnymi na brzegu.
Macierz pochodnej w punkcie brzegowym pozostaje
nieosobliwa. Otrzymujemy więc wspólną regularność $C^1$
odwrotności aż do końców przedziału czasu.

Dla dowolnego rozwiązania $u$ postawmy
$U(t;\xi)=u(t,X(t;\xi))$. Z \eqref{eq:pde2-chain}
wynika $\dot U+c(t,X(t;\xi))U=f(t,X(t;\xi))$.
Mnożąc przez czynnik całkujący $e^{C(t;\xi)}$,
otrzymujemy \eqref{eq:pde2-solution}.

Odwrotnie, prawa strona \eqref{eq:pde2-solution} wyznacza
funkcję $U\in C^1$; dzięki dyfeomorfizmowi $X_t$
definiuje ona $u(t,x)=U(t;X_t^{-1}(x))$ klasy $C^1$.
Różniczkowanie po charakterystykach daje
\eqref{eq:pde2-linear}. Każde inne rozwiązanie prowadzi
do tej samej funkcji $U$, co kończy dowód jednoznaczności.
\end{proof}

\begin{przyklad}[Transport w polu dylatacji]
Niech $b(x)=x$, $c=f=0$, $u_0\in C^1(\R^d)$.
Wówczas $X(t;\xi)=e^t\xi$ i
$u(t,x)=u_0(e^{-t}x)$. Profil rozszerza się
wraz z przepływem, ale wartości wzdłuż charakterystyk
nie zmieniają się.
\end{przyklad}

\begin{uwaga}[Odwracalność przepływu a odwracalność ewolucji]
W twierdzeniu \ref{thm:pde2-linear} w skończonym czasie
przepływ przestrzenny jest odwracalny. Nie oznacza to,
że każdy operator ewolucyjny PDE jest odwracalny
w dobrze postawionej klasie funkcji: dla równania ciepła
występuje wygładzanie, a dla transportu quasiliniowego
charakterystyki mogą przestać tworzyć układ współrzędnych.
\end{uwaga}

\section{Jak zmieniają się objętość i masa?}
\label{sec:pde2-jacobian}

Transport wartości skalarnej i transport gęstości to
dwa różne zagadnienia. Dla tego samego pola $b$ porównajmy
$u_t+b\cdot\nabla u=0$ oraz
$\rho_t+\diver_x(\rho b)=0$. Drugie równanie nazywamy
\emph{równaniem ciągłości}; w nim dodatkowo występuje
czynnik $(\diver_x b)\rho$.

\begin{stwierdzenie}[Wzór Jacobiego--Liouville'a]
\label{prop:pde2-liouville}
Niech $b$ spełnia założenia twierdzenia
\ref{thm:pde2-linear}, a $X$ będzie przepływem
\eqref{eq:pde2-flow-ode}. Dla $J(t;\xi)=\det D_\xi X(t;\xi)$
zachodzi
\begin{equation}\label{eq:pde2-jacobian}
 \frac{\mathrm d}{\mathrm dt}J(t;\xi)
 =\bigl(\diver_x b\bigr)(t,X(t;\xi))J(t;\xi),
 \qquad J(0;\xi)=1.
\end{equation}
W szczególności
$J(t;\xi)=\exp\!\bigl(\int_0^t
(\diver_x b)(s,X(s;\xi))\,\mathrm ds\bigr)>0$.
\end{stwierdzenie}
\begin{proof}
Ustalmy $\xi\in\R^d$ i oznaczmy
$A(t)=D_\xi X(t;\xi)$. Zapiszmy równanie
\eqref{eq:pde2-flow-ode} w postaci całkowej:
\[
 X(t;\xi)=\xi+\int_0^t b(s,X(s;\xi))\,\mathrm d\lambda_1(s).
\]
Na iloczynie $[0,T]$ i dowolnego zwartego otoczenia $\xi$
funkcja podcałkowa oraz jej pochodna względem $\xi$
są ciągłe, a więc ograniczone. Można zatem różniczkować
pod znakiem całki. Z reguły łańcuchowej wynika
\[
 A(t)=I_d+\int_0^t D_xb(s,X(s;\xi))A(s)\,\mathrm d\lambda_1(s),
\]
gdzie $I_d$ jest macierzą jednostkową rozmiaru $d$.
Prawa strona jest klasy $C^1$ względem $t$, ponieważ
funkcja podcałkowa jest ciągła. Twierdzenie o pochodnej
całki o zmiennej górnej granicy daje
$\dot A(t)=D_xb(t,X(t;\xi))A(t)$ oraz $A(0)=I_d$.
W końcach przedziału pochodne rozumiemy jednostronnie.
Macierz $A(t)$ jest odwracalna na mocy twierdzenia
\ref{thm:pde2-linear}. Zatem wzór na pochodną wyznacznika
daje $\frac{\mathrm d}{\mathrm dt}\det A
=\det A\,\tr(A^{-1}\dot A)$.
Cykliczność śladu i $\dot A=D_xb\,A$ prowadzą do
$\tr(A^{-1}\dot A)=\tr(D_xb)=\diver_xb$.
Oznaczmy $h(t)=(\diver_xb)(t,X(t;\xi))$ oraz
$H(t)=\int_0^t h(s)\,\mathrm d\lambda_1(s)$.
Z równania \eqref{eq:pde2-jacobian} wynika
$\frac{\mathrm d}{\mathrm dt}(e^{-H(t)}J(t;\xi))=0$.
Ponieważ $J(0;\xi)=\det I_d=1$ oraz $H(0)=0$,
otrzymujemy $J(t;\xi)=e^{H(t)}>0$.
\end{proof}

\begin{twierdzenie}[Transport gęstości i zachowanie masy]
\label{thm:pde2-mass}
Niech $b\in C^2([0,T]\times\R^d;\R^d)$ spełnia
warunki wzrostu i lipschitzowskości z twierdzenia
\ref{thm:pde2-linear}. Zakładamy ponadto, że wszystkie
pochodne przestrzenne $b$ rzędu $1$ i $2$ są ograniczone.
Dla dowolnej $\rho_0\in C_c^1(\R^d)$ istnieje dokładnie jedno
rozwiązanie $\rho\in C^1([0,T]\times\R^d)$ równania
\[
 \rho_t+\diver_x(\rho b)=0,\qquad \rho(0,x)=\rho_0(x).
\]
Jest ono zadane zależnością
\begin{equation}\label{eq:pde2-mass-flow}
 \rho(t,X(t;\xi))J(t;\xi)=\rho_0(\xi).
\end{equation}
Ponadto $\int_{\R^d}\rho(t,x)\,\mathrm d\lambda_d(x)
=\int_{\R^d}\rho_0(x)\,\mathrm d\lambda_d(x)$ dla $t\in[0,T]$.
\end{twierdzenie}
\begin{proof}
Ponieważ $\diver_x(\rho b)=b\cdot\nabla\rho+
(\diver_xb)\rho$, twierdzenie \ref{thm:pde2-linear}
można zastosować z $c=\diver_xb$, $f=0$.
Założenie $b\in C^2$ zapewnia, że $\diver_x b\in C^1$,
więc twierdzenie zapewnia istnienie i jednoznaczność
rozwiązania klasy $C^1$. Wzdłuż charakterystyki zachodzi
$\frac{\mathrm d}{\mathrm dt}\rho(t,X(t;\xi))
=-(\diver_x b)(t,X(t;\xi))\rho(t,X(t;\xi))$.
Po zastosowaniu reguły iloczynu i \eqref{eq:pde2-jacobian}
otrzymujemy
\[
 \frac{\mathrm d}{\mathrm dt}
 \bigl(\rho(t,X(t;\xi))J(t;\xi)\bigr)=0.
\]
W chwili $t=0$ iloczyn wynosi $\rho_0(\xi)$,
co dowodzi \eqref{eq:pde2-mass-flow}.
Ponieważ $J(t;\xi)>0$, funkcja $\rho(t,X_t(\xi))$
jest niezerowa wtedy i tylko wtedy, gdy $\rho_0(\xi)\ne0$.
Homeomorfizm $X_t$ zachowuje domknięcia zbiorów,
więc $\supp\rho(t,\cdot)=X_t(\supp\rho_0)$.
Jest to zbiór zwarty jako ciągły obraz zbioru zwartego.
Obie gęstości są ciągłe i mają zwarte nośniki,
a zatem są całkowalne również wtedy, gdy zmieniają znak.
Możemy więc zastosować twierdzenie o zmianie zmiennych
dla dyfeomorfizmu $x=X_t(\xi)$.
Dodatniość $J$ i \eqref{eq:pde2-mass-flow} dają
\[
 \int_{\R^d}\rho(t,x)\,\mathrm d\lambda_d(x)
 =\int_{\R^d}\rho(t,X(t;\xi))J(t;\xi)
 \,\mathrm d\lambda_d(\xi)
 =\int_{\R^d}\rho_0(\xi)\,\mathrm d\lambda_d(\xi).
\]
\end{proof}

\begin{uwaga}[Warunki regularności]
Wzór \eqref{eq:pde2-mass-flow} zachowuje sens również
przy słabszych założeniach. W poprzednim dowodzie
zastosowaliśmy dodatkową regularność $b$, aby pozostać
w klasie rozwiązań klasycznych; teoria dystrybucyjna
pozwoli te wymagania osłabić.
\end{uwaga}

\begin{stwierdzenie}[Kontrola jakobianu i norm przy transporcie]
\label{prop:pde-v31-transport-lp}
Niech $T>0$, $d\ge1$, a $b$ spełnia założenia
twierdzenia~\ref{thm:pde2-linear} z $\|D_xb\|\le M$.
Niech $X_t$ i $J(t;\xi)=\det D_\xi X_t(\xi)$ będą
odpowiednio przepływem i jego jakobianem. Wówczas
\[
 e^{-dMt}\le J(t;\xi)\le e^{dMt}
 \qquad(0\le t\le T,\ \xi\in\R^d).
\]
Jeśli $u_0\in C^1(\R^d)\cap L^p(\R^d)$ dla
$1\le p<\infty$, a $u$ rozwiązuje
$u_t+b\cdot\nabla u=0$, $u(0)=u_0$, to
\begin{equation}\label{eq:pde-v31-transport-lp}
 e^{-dMt/p}\|u_0\|_p\le\|u(t,\cdot)\|_p
 \le e^{dMt/p}\|u_0\|_p.
\end{equation}
Dla ograniczonego $u_0\in C^1(\R^d)$ zachodzi ponadto
$\|u(t,\cdot)\|_\infty=\|u_0\|_\infty$.
Jeśli $\diver_xb=0$, to $J\equiv1$ i wszystkie wymienione
normy $L^p$ są zachowane.
\end{stwierdzenie}
\begin{proof}
Niech $(e_1,\ldots,e_d)$ będzie standardową bazą
ortonormalną $\R^d$. Dla każdej rzeczywistej macierzy
$A$ rozmiaru $d$, z nierówności Cauchy'ego--Schwarza,
\[
 |\tr A|\le\sum_{j=1}^d|\langle Ae_j,e_j\rangle|
 \le\sum_{j=1}^d |Ae_j|\le d\|A\|.
\]
Stosując to oszacowanie do $A=D_xb(t,x)$, otrzymujemy
$|(\diver_xb)(t,x)|\le dM$.
Wzór~\eqref{eq:pde2-jacobian} i jego rozwiązanie dają
pierwsze oszacowanie. Wzdłuż charakterystyk
$u(t,X_t(\xi))=u_0(\xi)$. Ponieważ $X_t$ jest
dyfeomorfizmem, wzór na zmianę zmiennych daje
\[
 \|u(t,\cdot)\|_p^p
 =\int_{\R^d}|u_0(\xi)|^pJ(t;\xi)\,
   \mathrm d\lambda_d(\xi).
\]
Z oszacowań na $J$ wynika
$e^{-dMt}\|u_0\|_p^p\le\|u(t,\cdot)\|_p^p
\le e^{dMt}\|u_0\|_p^p$. W szczególności całka jest
skończona; wyciągnięcie pierwiastka stopnia $p$
daje \eqref{eq:pde-v31-transport-lp}.
Dla ograniczonej funkcji ciągłej supremum modułu jest
równe supremum istotnemu względem miary Lebesgue'a:
jeśli liczba jest mniejsza od supremum, zbiór punktów,
w których moduł ją przekracza, jest niepusty i otwarty,
a więc ma dodatnią miarę. Ponieważ $X_t$ jest bijekcją,
$\sup_x|u(t,x)|=\sup_\xi|u_0(\xi)|$; stąd wynika równość
norm $L^\infty$. W przypadku zerowej dywergencji
równanie na $J$ ma rozwiązanie stałe równe $1$.
\end{proof}

\section{Niecharakterystyczne dane Cauchy'ego na hiperpowierzchni}
\label{sec:pde2-noncharacteristic}

Rozważmy równanie bez wyróżnionego czasu
\begin{equation}\label{eq:pde2-static}
 a(x)\cdot\nabla u(x)+c(x)u(x)=f(x),\qquad x\in\Omega,
\end{equation}
gdzie $\Omega\subset\R^n$ jest otwarta, $n\geq2$.
Gdy dane zadano na hiperpowierzchni $\Sigma$, istotne
jest, czy charakterystyki przecinają ją poprzecznie.

\begin{definicja}[Niecharakterystyczność dla transportu]
Niech $\Omega\subset\R^n$ będzie otwarta, $a\in C^1(\Omega;\R^n)$,
a $\Sigma\subset\Omega$ hiperpowierzchnią klasy $C^2$.
Mówimy, że $\Sigma$ jest niecharakterystyczna w punkcie
$x_0\in\Sigma$ dla operatora $a\cdot\nabla$, jeżeli
$a(x_0)\notin T_{x_0}\Sigma$. Jeśli w pewnym otwartym
otoczeniu $O\subset\Omega$ punktu $x_0$ mamy
$\Sigma\cap O=\{x\in O:\phi(x)=0\}$, gdzie $\phi\in C^2(O)$
i $\mathrm d\phi(x_0)\ne0$, to
$T_{x_0}\Sigma=\ker\mathrm d\phi(x_0)$. Niecharakterystyczność
jest zatem równoważna warunkowi
$a(x_0)\cdot\nabla\phi(x_0)\ne0$.
\end{definicja}

\begin{twierdzenie}[Lokalne zagadnienie Cauchy'ego]
\label{thm:pde2-cauchy}
Niech $\Omega\subset\R^n$ będzie otwarta,
$a\in C^1(\Omega;\R^n)$, $c,f\in C^1(\Omega)$,
$\Sigma\subset\Omega$
hiperpowierzchnią klasy $C^2$, a $x_0\in\Sigma$.
Niech $g\in C^1(\Sigma)$ w otoczeniu $x_0$.
Jeśli $\Sigma$ jest niecharakterystyczna w $x_0$ dla
operatora $a\cdot\nabla$, to istnieje otwarte otoczenie
$U\subset\Omega$ punktu $x_0$, na którym zagadnienie
\eqref{eq:pde2-static} z danymi $u|_{U\cap\Sigma}=g$
ma rozwiązanie $u\in C^1(U)$. Otoczenie $U$ można wybrać
tak, aby rozwiązanie było jedyne na $U$.
W szczególności rozwiązanie jest jednoznaczne jako zarodek
w $x_0$: dwa rozwiązania lokalne są równe na pewnym
wspólnym otoczeniu tego punktu.
\end{twierdzenie}
\begin{proof}
Wybierzmy lokalną parametryzację
$\gamma\colon V\to\Sigma$, gdzie $V\subset\R^{n-1}$
jest otwartym otoczeniem zera, $\gamma$ jest klasy $C^2$,
$\gamma(0)=x_0$, a wektory
$\partial_{y_1}\gamma(0),\ldots,\partial_{y_{n-1}}\gamma(0)$
są liniowo niezależne. Niech $\Phi_s(x)$ oznacza lokalny
przepływ pola $a$, tj. $\partial_s\Phi_s(x)=a(\Phi_s(x))$
i $\Phi_0(x)=x$. Twierdzenie o lokalnym istnieniu
przepływu pola klasy $C^1$ zapewnia wspólny przedział
$(-\delta_0,\delta_0)$ dla punktów początkowych w pewnym
otoczeniu $x_0$ oraz regularność $C^1$ względem czasu
i punktu początkowego. Zmniejszając $V$, możemy przyjąć,
że $\gamma(V)$ leży w tym otoczeniu oraz w dziedzinie danych $g$.
Zdefiniujmy $\Psi(s,y)=\Phi_s(\gamma(y))$ na
$(-\delta_0,\delta_0)\times V$.

Macierz pochodnej $D\Psi(0,0)$ ma kolumny
$a(x_0)$ oraz $\partial_{y_j}\gamma(0)$. Założenie
niecharakterystyczności mówi dokładnie, że tworzą bazę
$\R^n$. Z twierdzenia o funkcji odwrotnej istnieją
otoczenia $W$ punktu $(0,0)$ i $U$ punktu $x_0$,
dla których $\Psi\colon W\to U$ jest dyfeomorfizmem.
Wybierzmy lokalną funkcję definiującą $\phi$ dla $\Sigma$
w otoczeniu $x_0$, jak w poprzedniej definicji.
Ponieważ
$\partial_s(\phi\circ\Psi)(0,0)=\mathrm d\phi(x_0)a(x_0)\ne0$,
po zmniejszeniu otoczenia pochodna ta ma stały niezerowy znak.
Wybierzmy iloczyn $W'=(-\delta,\delta)\times V'\subset W$
zawierający $(0,0)$ tak, aby $\Psi(W')$ leżał w dziedzinie
$\phi$ i aby powyższy warunek znaku zachodził na $W'$.
Dla każdego $y\in V'$ mamy $\phi(\Psi(0,y))=0$.
Ścisła monotoniczność funkcji $s\mapsto\phi(\Psi(s,y))$
pokazuje, że jej jedynym zerem jest $s=0$.
Przyjmując $U=\Psi(W')$, otrzymujemy otwarte otoczenie $x_0$,
na którym $\Psi|_{W'}$ jest dyfeomorfizmem i
$U\cap\Sigma=\Psi(\{0\}\times V')$.
W szczególności cały odcinek od $(0,y)$ do $(s,y)$
leży w $W'$ dla każdego $(s,y)\in W'$.

Niech $v(s,y)=u(\Psi(s,y))$. Równanie
\eqref{eq:pde2-static} sprowadza się do ODE
$\partial_sv+c(\Psi(s,y))v=f(\Psi(s,y))$
z danymi $v(0,y)=g(\gamma(y))$. Dla
$C(s,y)=\int_0^s c(\Psi(r,y))\,\mathrm dr$ rozwiązanie ma
postać
\[
 v(s,y)=e^{-C(s,y)}g(\gamma(y))+
 \int_0^s e^{-(C(s,y)-C(r,y))}
 f(\Psi(r,y))\,\mathrm dr.
\]
Całki dla $s<0$ rozumiemy jako całki z odwróconą
orientacją granic. Na zwartych podzbiorach iloczynu $W'$
współczynniki $c\circ\Psi$, $f\circ\Psi$ oraz ich pierwsze
pochodne są ciągłe i lokalnie ograniczone. Reguła Leibniza
dla całek z parametrem, zastosowana na zwartych otoczeniach
zawierających odpowiednie odcinki całkowania, daje $v\in C^1(W')$.
Różniczkowanie wzoru względem $s$ daje
$\partial_s v+c(\Psi(s,y))v=f(\Psi(s,y))$, a podstawienie
$s=0$ daje $v(0,y)=g(\gamma(y))$.
Złożenie $u=v\circ(\Psi|_{W'})^{-1}$ jest klasy $C^1$.
Ponieważ $\partial_s\Psi=a\circ\Psi$, reguła łańcuchowa
pokazuje, że $u$ spełnia \eqref{eq:pde2-static} na $U$.
Równość $U\cap\Sigma=\Psi(\{0\}\times V')$ zapewnia
spełnienie danych na całym $U\cap\Sigma$.
Każde inne rozwiązanie na $U$ po złożeniu z $\Psi$
spełnia ten sam skalarny problem ODE na $(-\delta,\delta)$
dla każdego $y\in V'$, więc z jednoznaczności ODE jest równe $u$.
Dla dwóch rozwiązań zdefiniowanych na innych otoczeniach
wybieramy mniejszy iloczyn zawarty w przeciwobrazie
ich wspólnej dziedziny przez $\Psi$. Ten sam argument
dowodzi jednoznaczności zarodka.
\end{proof}

\begin{center}
\begin{tikzpicture}[>=Stealth,scale=.92]
 \draw[->] (-3.3,-1.35)--(3.8,-1.35) node[right] {$x_1$};
 \draw[->] (-3.1,-1.55)--(-3.1,2.5) node[above] {$x_2$};
 \draw[very thick] (-2.55,-.65)..controls (-.8,-.8) and (1.3,-.18)..(2.9,-.45);
 \node[below] at (2.6,-.6) {$\Sigma$};
 \foreach \x/\yy in {-1.75/-.64,-.65/-.52,.55/-.37,1.65/-.34} {
   \draw[blue!65!black,thick,->] (\x,\yy)--({\x+.48},{\yy+1.85});
   \fill (\x,\yy) circle (1.5pt);
 }
 \fill[red!70!black] (.55,-.37) circle (2.4pt);
 \node[below] at (.55,-.76) {$x_0$};
 \node[above] at (1.7,1.4) {charakterystyki};
\end{tikzpicture}
\par\smallskip\small
Dane są zadawane na $\Sigma$; przechodzące przez nią
poprzecznie charakterystyki tworzą lokalny układ współrzędnych.
\end{center}

\begin{kontrprzyklad}[Dane na powierzchni charakterystycznej]
\label{ex:pde2-char-data}
Rozważmy $u_x=0$ na $\R^2$, gdzie $u=u(x,y)$,
oraz dane $u(x,0)=g(x)$ na prostej $\Sigma=\{y=0\}$.
Każde rozwiązanie ma postać $u(x,y)=h(y)$, stąd
istnienie wymaga, by $g$ była stała. Gdy $g\equiv c$,
rozwiązań jest nieskończenie wiele: można wybrać
dowolne $h\in C^1(\R)$ z $h(0)=c$. Prosta $y=0$
jest charakterystyczna, bo pole $(1,0)$ leży
w jej przestrzeni stycznej.
\end{kontrprzyklad}

\section{Quasiliniowe równania pierwszego rzędu}
\label{sec:pde2-quasilinear}

Teraz prędkość charakterystyki może zależeć od niewiadomej:
\begin{equation}\label{eq:pde2-quasilinear}
 u_t+a(t,x,u)\cdot\nabla_xu=b(t,x,u),
 \qquad u(0,x)=u_0(x).
\end{equation}
W odróżnieniu od przypadku liniowego nie znamy przepływu
z góry: równocześnie musimy wyznaczać jego położenie
$X$ oraz wartość $U$ rozwiązania na charakterystyce.

\begin{twierdzenie}[Lokalna konstrukcja i jednoznaczność]
\label{thm:pde2-quasilinear}
Niech $a\colon(-T,T)\times\R^d\times\R\to\R^d$
i $b\colon(-T,T)\times\R^d\times\R\to\R$
będą klasy $C^2$, a $u_0\in C^2(\R^d)$.
Dla dowolnego $\xi_0\in\R^d$ istnieją otwarte otoczenie
$V\subset\R^d$ punktu $\xi_0$, liczba $0<\delta<T$
i otwarty zbiór $Q\subset(-\delta,\delta)\times\R^d$
zawierający $(0,\xi_0)$, na którym zagadnienie
\eqref{eq:pde2-quasilinear} ma rozwiązanie $u\in C^1(Q)$
z danymi $u(0,x)=u_0(x)$ dla $(0,x)\in Q$.
Rozwiązanie jest jednoznaczne jako zarodek klasy $C^1$
w otoczeniu $(0,\xi_0)$.

Dokładniej, niech $(X,U)$ rozwiązują układ ODE
\begin{equation}\label{eq:pde2-quasi-chars}
 \begin{aligned}
 \dot X(t;\xi)&=a(t,X(t;\xi),U(t;\xi)),& X(0;\xi)&=\xi,\\
 \dot U(t;\xi)&=b(t,X(t;\xi),U(t;\xi)),& U(0;\xi)&=u_0(\xi).
 \end{aligned}
\end{equation}
Układ ma rozwiązania dla wszystkich $(t,\xi)\in
(-\delta,\delta)\times V$, a po odpowiednim wyborze
$\delta,V$ odwzorowanie
$(t,\xi)\mapsto(t,X(t;\xi))$ jest dyfeomorfizmem
tego iloczynu na $Q$. Rozwiązanie spełnia
$u(t,X(t;\xi))=U(t;\xi)$.
\end{twierdzenie}
\begin{proof}
Prawa strona układu \eqref{eq:pde2-quasi-chars}
jest klasy $C^2$ w otwartym zbiorze
$(-T,T)\times\R^{d+1}$, a odwzorowanie danych
$\xi\mapsto(\xi,u_0(\xi))$ jest klasy $C^2$.
Lokalne twierdzenie o istnieniu ODE i zależności od danych
zapewnia wspólny przedział istnienia wokół zera dla
punktów początkowych z pewnego otoczenia
$(\xi_0,u_0(\xi_0))$. Z ciągłości odwzorowania danych
wynika więc istnienie rozwiązań $(X,U)$ klasy co najmniej
$C^1$ względem $(t,\xi)$ na iloczynie otwartego przedziału
i otoczenia $\xi_0$.
Dla $\Theta(t,\xi)=(t,X(t;\xi))$ mamy
\[
 D\Theta(0,\xi_0)=
 \begin{pmatrix}1&0\\
 a(0,\xi_0,u_0(\xi_0))&I_d
 \end{pmatrix},
 \qquad \det D\Theta(0,\xi_0)=1.
\]
Stosujemy twierdzenie o funkcji odwrotnej. W otoczeniu,
na którym $\Theta$ jest dyfeomorfizmem, wybieramy mniejszy
iloczyn $(-\delta,\delta)\times V$ zawierający $(0,\xi_0)$,
z $\delta<T$, i kładziemy $Q=\Theta((-\delta,\delta)\times V)$.
Zbiór $Q$ jest otwarty, a ponieważ $\Theta$ zachowuje
pierwszą współrzędną, leży w $(-\delta,\delta)\times\R^d$.

Niech $\xi=\Xi(t,x)$ oznacza jej lokalną odwrotność
współrzędnych przestrzennych. Kładziemy
$u(t,x)=U(t;\Xi(t,x))$. Różniczkując tożsamość
$u(t,X(t;\xi))=U(t;\xi)$ po $t$ przy ustalonym $\xi$,
otrzymujemy
\[
 u_t(t,X)+\nabla_xu(t,X)\cdot a(t,X,U)
 =b(t,X,U).
\]
Ponieważ $U=u(t,X)$, jest to \eqref{eq:pde2-quasilinear}.
W chwili $t=0$ mamy $X(0;\xi)=\xi$ i
$U(0;\xi)=u_0(\xi)$, więc warunek początkowy zachodzi.

Jeśli $v\in C^1$ jest innym lokalnym rozwiązaniem z tymi
danymi na otwartym otoczeniu $(0,\xi_0)$, pole
$(t,y)\mapsto a(t,y,v(t,y))$ jest tam klasy $C^1$.
Lokalne twierdzenie o istnieniu ODE daje wspólny przedział
$(-\varepsilon,\varepsilon)$ dla rozwiązań
$\dot Y(t;\xi)=a(t,Y(t;\xi),v(t,Y(t;\xi)))$,
$Y(0;\xi)=\xi$, dla $\xi$ z pewnego otoczenia $V_1\subset V$
punktu $\xi_0$. Wybieramy $\varepsilon\le\delta$ i $V_1$
tak, aby trajektorie pozostawały w dziedzinie $v$,
a $(0,\xi)$ należało do zbioru danych początkowych.
Funkcja $Z(t;\xi)=v(t,Y(t;\xi))$ spełnia z reguły
łańcuchowej i równania dla $v$
$\partial_t Z=b(t,Y,Z)$ oraz $Z(0;\xi)=u_0(\xi)$.
Para $(Y,Z)$ rozwiązuje więc \eqref{eq:pde2-quasi-chars}
z tymi samymi danymi co $(X,U)$. Jednoznaczność ODE daje
$(Y,Z)=(X,U)$ na całym wspólnym iloczynie.
W konsekwencji $v(t,X(t;\xi))=U(t;\xi)=u(t,X(t;\xi))$.
Obraz $(-\varepsilon,\varepsilon)\times V_1$ przez
dyfeomorfizm $\Theta$ jest otwartym otoczeniem $(0,\xi_0)$,
na którym $v=u$.
\end{proof}

\begin{uwaga}[Co może uniemożliwić przedłużenie?]
Lokalna konstrukcja wymaga zarówno istnienia rozwiązań
układu ODE, jak i odwracalności odwzorowania $\xi\mapsto X(t;\xi)$.
W zagadnieniach quasiliniowych nawet przy globalnym
istnieniu charakterystyk może wystąpić
$\det D_\xi X=0$: wtedy klasyczne współrzędne
charakterystyczne zawodzą. W modelach skalarnych jest
to typowy mechanizm utraty regularności rozwiązania.
\end{uwaga}

\section{Fale proste i kryterium wybuchu pochodnej}
\label{sec:pde2-blowup}

Rozważmy w jednym wymiarze równanie autonomiczne
\begin{equation}\label{eq:pde2-simplewave}
 u_t+a(u)u_x=0,\qquad u(0,x)=u_0(x),
\end{equation}
gdzie $a\in C^2(\R)$, $u_0\in C^2(\R)$.
Wzdłuż charakterystyki wartość $u$ jest stała, więc
\begin{equation}\label{eq:pde2-simplewave-map}
 X(t;\xi)=\xi+t\,a(u_0(\xi)),\qquad
 U(t;\xi)=u_0(\xi).
\end{equation}
Różniczkując względem $\xi$, dostajemy
$X_\xi(t;\xi)=1+t\,a'(u_0(\xi))u_0'(\xi)$.

\begin{twierdzenie}[Pierwsza osobliwość gradientu]
\label{thm:pde2-breaking}
Niech $a,u_0\in C^2(\R)$ i załóżmy, że
$a\circ u_0$ jest ograniczona, a funkcja
$q(\xi)=a'(u_0(\xi))u_0'(\xi)$ jest ograniczona z dołu.
Jeżeli $m=\inf_{\xi\in\R}q(\xi)<0$, określmy
$T_*=-1/m$. Wtedy dla każdego $0\leq t<T_*$
odwzorowanie $X_t(\xi)=\xi+t a(u_0(\xi))$
jest ściśle rosnącym dyfeomorfizmem $\R\to\R$ i
\eqref{eq:pde2-simplewave-map} definiuje rozwiązanie
klasy $C^1$ na $[0,T_*)\times\R$. Jego pochodna wynosi
\begin{equation}\label{eq:pde2-gradient-formula}
 u_x(t,X(t;\xi))=
 \frac{u_0'(\xi)}{1+t a'(u_0(\xi))u_0'(\xi)}.
\end{equation}
Jeśli ponadto minimum $m$ jest osiągane w punkcie
$\xi_*$, to $|u_x(t,X(t;\xi_*))|\to\infty$
gdy $t\uparrow T_*$. Nie istnieje zatem przedłużenie
rozwiązania jako funkcji $C^1$ na otoczenie
$(T_*,X(T_*;\xi_*))$.
\end{twierdzenie}
\begin{proof}
Dla $0\leq t<T_*$ mamy
$X_\xi(t;\xi)=1+tq(\xi)\geq1+tm>0$.
Odwzorowanie $X_t$ jest więc ściśle rosnące.
Z ograniczoności $a\circ u_0$ wynika
$X_t(\xi)-\xi=O(1)$ dla $\xi\to\pm\infty$,
więc $X_t(\xi)\to\pm\infty$ przy $\xi\to\pm\infty$.
Z ciągłości wynika surjektywność, a ze ścisłej
monotoniczności injektywność $X_t$.
Odwzorowanie $(t,\xi)\mapsto(t,X(t;\xi))$ ma wyznacznik
macierzy pochodnej równy $1+tq(\xi)>0$.
Twierdzenie o funkcji odwrotnej daje lokalne odwrotności
klasy $C^1$; ich zgodność wynika z bijektywności $X_t$.
Przy $t=0$ stosujemy to twierdzenie do tego samego wzoru
na $X(t;\xi)$, określonego także dla ujemnych $t$.
W ten sposób $\Xi(t,x)=X_t^{-1}(x)$ jest klasy $C^1$
względem obu zmiennych aż do $t=0$.

Definiujemy $u(t,x)=u_0(\Xi(t,x))$. Jest to funkcja klasy
$C^1$, a ponieważ $X_0$ jest identycznością, mamy
$u(0,x)=u_0(x)$. Z równości $u(t,X(t;\xi))=u_0(\xi)$
przez różniczkowanie względem $t$ otrzymujemy
\[
 u_t(t,X(t;\xi))+u_x(t,X(t;\xi))a(u_0(\xi))=0.
\]
Ponieważ $u_0(\xi)=u(t,X(t;\xi))$ i $X_t$ jest surjekcją,
jest to równanie \eqref{eq:pde2-simplewave} w każdym punkcie
dziedziny. Różniczkowanie tej samej tożsamości względem
$\xi$ daje $u_x(t,X(t;\xi))X_\xi(t;\xi)=u_0'(\xi)$.
Dzieląc przez dodatnią pochodną $X_\xi$, uzyskujemy
\eqref{eq:pde2-gradient-formula}.

Jeżeli $q(\xi_*)=m<0$, to
$a'(u_0(\xi_*))u_0'(\xi_*)\ne0$; w szczególności
$u_0'(\xi_*)\ne0$. Mianownik
\eqref{eq:pde2-gradient-formula} przy $\xi_*$
dąży do zera od strony dodatniej, a licznik jest niezerowy.
Ponadto $X(t;\xi_*)\to X(T_*;\xi_*)$, ponieważ wzór
na charakterystykę jest ciągły także w $t=T_*$.
Zatem moduł pochodnej dąży do nieskończoności wzdłuż
krzywej dochodzącej do wskazanego punktu. Ciągłość
pochodnej przestrzennej ewentualnego przedłużenia
w tym punkcie przeczyłaby tej granicy.
\end{proof}

\begin{przyklad}[Burgers: dokładny czas powstania osobliwości]
\label{ex:pde2-burgers}
Dla równania $u_t+u u_x=0$ mamy $a(u)=u$.
Wybierzmy $u_0(x)=-\tanh x$. Wówczas
$q(x)=u_0'(x)=-1/\cosh^2x$ osiąga minimum $-1$
w $x=0$, zaś $a\circ u_0$ jest ograniczona.
Twierdzenie \ref{thm:pde2-breaking} daje $T_*=1$,
$X(t;\xi)=\xi-t\tanh\xi$ oraz
$u_x(t,0)=-1/(1-t)$ dla $t<1$.
Charakterystyki istnieją jako krzywe gładkie także
później, ale nie mogą już wyznaczać globalnego
klasycznego rozwiązania w tym samym schemacie.
\end{przyklad}

\begin{center}
\begin{tikzpicture}[>=Stealth,scale=.86]
 \draw[->] (-3.4,0)--(3.6,0) node[right] {$x$};
 \draw[->] (-3.1,-.12)--(-3.1,3.75) node[above] {$t$};
 \foreach \xi in {-2.4,-1.6,-.8,-.4,.4,.8,1.6,2.4} {
  \pgfmathsetmacro{\dd}{-(exp(2*\xi)-1)/(exp(2*\xi)+1)}
  \draw[blue!65!black,thick] (\xi,0)--({\xi+0.95*\dd},2.95);
 }
 \draw[densely dashed] (-3.1,3.1)--(3.35,3.1);
 \node[anchor=west] at (3.55,3.08) {$t\approx T_*$};
 \node[below] at (0,0) {$0$};
 \node[align=left,anchor=west] at (3.6,1.45) {zbliżanie się\\charakterystyk};
\end{tikzpicture}
\par\smallskip\small
Dla malejących danych Burgersa charakterystyki
zbliżają się do siebie. Rysunek ma charakter schematyczny;
czas i nachylenia przeskalowano niezależnie.
\end{center}

\begin{uwaga}[Eksperyment obliczeniowy]
W rozdziale~\ref{chap:pde-code} zamieszczono kod Python,
który dla przykładu Burgersa z danymi $u_0=-\tanh x$
rysuje charakterystyki i profile rozwiązań dla
$t\in\{0,0.5,0.9,0.99\}$. Program korzysta z odwracania
monotonicznego odwzorowania $\xi\mapsto X(t;\xi)$ przez
interpolację. Porównanie z dokładnym wzorem
\eqref{eq:pde2-gradient-formula} pozwala sprawdzić,
że $u_x(t,0)=-1/(1-t)$, podczas gdy $|u(t,x)|\leq1$.
Nie jest to algorytm kontynuacji rozwiązania przez szok:
dla $t\geq1$ odwzorowanie etykiet nie jest już globalnie
odwracalna w klasie $C^1$.
\end{uwaga}

\begin{uwaga}[Znaczenie czasu $T_*$]
Twierdzenie \ref{thm:pde2-breaking} dotyczy dokładnie
zadanego równania $u_t+a(u)u_x=0$, bez źródła,
i wskazanej klasy danych. Nie jest ogólnym twierdzeniem
o czasie istnienia dowolnych równań quasiliniowych.
Gdy $q(\xi)\geq0$ dla wszystkich $\xi$, charakterystyki
nie przecinają się dla $t\geq0$, o ile
$a\circ u_0$ pozostaje ograniczona; w takim przypadku
wzór daje globalne rozwiązanie klasy $C^1$.
\end{uwaga}

\section{Od przecinających się charakterystyk do fal uderzeniowych}
\label{sec:pde2-shocks}

Równanie $u_t+a(u)u_x=0$ można, gdy $a=f'$,
zapisać w postaci prawa zachowania
\begin{equation}\label{eq:pde2-conservation}
 u_t+\partial_x f(u)=0.
\end{equation}
Po pojawieniu się nieciągłości wzór z pochodnymi
klasycznymi przestaje mieć sens. Potrzebne jest
sformułowanie całkowe, które później zinterpretujemy
w języku dystrybucji.

\begin{definicja}[Rozwiązanie słabe prawa zachowania]
\label{def:pde2-weak}
Niech $f\in C^1(\R)$, $u_0\in L^1_{\mathrm{loc}}(\R)$,
a $u\in L^1_{\mathrm{loc}}([0,T)\times\R)$, przy czym
$f(u)$ jest lokalnie całkowalna. Funkcję $u$ nazywamy
\emph{rozwiązaniem słabym} zagadnienia
\eqref{eq:pde2-conservation} z danymi $u_0$, jeżeli
dla każdej $\varphi\in C_c^\infty([0,T)\times\R)$
zachodzi
\begin{equation}\label{eq:pde2-weak-form}
 \int_0^T\!\int_\R
 \bigl(u(t,x)\varphi_t(t,x)+f(u(t,x))\varphi_x(t,x)\bigr)
 \,\mathrm d\lambda_1(x)\,\mathrm dt
 +\int_\R u_0(x)\varphi(0,x)\,\mathrm d\lambda_1(x)=0.
\end{equation}
Przestrzeń $C_c^\infty([0,T)\times\R)$ oznacza
restrykcje funkcji gładkich w pewnym otoczeniu
$[0,T)\times\R$, których nośniki są zwarte i nie dochodzą
do czasu $T$.
\end{definicja}

\begin{stwierdzenie}[Warunek Rankine'a--Hugoniota]
\label{prop:pde2-rh}
Niech $f\in C^1(\R)$ i $s\in C^1((0,T))$.
Załóżmy, że $u$ jest klasy $C^1$ po obu stronach
krzywej $x=s(t)$, z jednostronnymi przedłużeniami
klasy $C^1$ do krzywej, i posiada ślady
$u_L(t)$ i $u_R(t)$ na tej krzywej oraz spełnia
\eqref{eq:pde2-conservation} klasycznie poza nią.
Wówczas $u$ spełnia równanie w sensie dystrybucji
na $(0,T)\times\R$ (bez warunku początkowego), tj.
$\int_0^T\int_\R (u\varphi_t+f(u)\varphi_x)\,
\mathrm d\lambda_1(x)\,\mathrm d\lambda_1(t)=0$
dla każdej $\varphi\in C_c^\infty((0,T)\times\R)$,
wtedy i tylko wtedy, gdy
\begin{equation}\label{eq:pde2-rh}
 s'(t)\bigl(u_R(t)-u_L(t)\bigr)
 =f(u_R(t))-f(u_L(t))
\end{equation}
dla każdego $t\in(0,T)$, przy ciągłości śladów.
\end{stwierdzenie}
\begin{proof}
Niech $\varphi\in C_c^\infty((0,T)\times\R)$.
Wybierzmy $0<a<b<T$ oraz $R>0$, tak aby nośnik
$\varphi$ leżał w $(a,b)\times(-R,R)$ oraz
$|s(t)|<R$ dla $t\in[a,b]$. Oznaczmy
$D_L=\{(t,x):a<t<b,\ -R<x<s(t)\}$ oraz
$D_R=\{(t,x):a<t<b,\ s(t)<x<R\}$.
Na każdym z tych zbiorów, z równania klasycznego
i reguły iloczynu, mamy
\[
 \partial_t(u\varphi)+\partial_x(f(u)\varphi)
 =u\varphi_t+f(u)\varphi_x.
\]
Całkujemy po obu zbiorach, stosując twierdzenie Gaussa
do pola $(u\varphi,f(u)\varphi)$ w zmiennych $(t,x)$.
Brzegi są kawałkami klasy $C^1$, pole ma jednostronne
przedłużenia klasy $C^1$, a $\varphi$ znika na zewnętrznych
krawędziach prostokąta. Na krzywej $(t,s(t))$
wektor normalny zewnętrzny dla $D_L$ ma kierunek
$(-s'(t),1)$, a dla $D_R$ kierunek przeciwny.
Element długości krzywej wynosi
$\sqrt{1+s'(t)^2}\,\mathrm d\lambda_1(t)$,
co skraca czynnik normalizujący normalnej.
Wkład od lewej strony wynosi więc
$\int_a^b(f(u_L)-s'u_L)\varphi(t,s(t))\,\mathrm d\lambda_1(t)$,
a od prawej strony
$\int_a^b(s'u_R-f(u_R))\varphi(t,s(t))\,\mathrm d\lambda_1(t)$.
Ich suma wynosi
\[
 \int_0^T\bigl(s'(t)[u](t)-[f(u)](t)\bigr)
 \varphi(t,s(t))\,\mathrm dt,
\]
gdzie $[u]=u_R-u_L$ oraz $[f(u)]=f(u_R)-f(u_L)$.
Ślad funkcji testowej na krzywej można zadawać dowolnie
bez zakładania gładkości samej krzywej ponad klasę $C^1$.
Jeśli zachodzi \eqref{eq:pde2-rh}, suma jest zerowa
dla każdej funkcji testowej, co dowodzi wystarczalności.
Dla dowodu konieczności załóżmy, że równanie zachodzi
w sensie dystrybucji. Niech $\eta\in C_c^\infty((0,T))$
i niech $J\subset(0,T)$ będzie zwartym przedziałem
zawierającym jej nośnik. Zwartość $s(J)$
pozwala wybrać $\chi\in C_c^\infty(\R)$, równą $1$
w otoczeniu $s(J)$. Wtedy
$\varphi(t,x)=\eta(t)\chi(x)$ jest klasy $C_c^\infty$
i spełnia $\varphi(t,s(t))=\eta(t)$ dla $t\in J$.
Zatem dla ciągłej funkcji
$b(t)=s'(t)[u](t)-[f(u)](t)$ zachodzi
$\int_0^T b(t)\eta(t)\,\mathrm dt=0$ dla każdego takiego $\eta$.
Gdyby $b(t_0)\ne0$, ciągłość dawałaby przedział, na którym
$b$ ma stały, niezerowy znak; test nieujemny o nośniku w tym
przedziale prowadziłby do sprzeczności. Stąd $b=0$
w każdym punkcie $(0,T)$.
Warunek rozstrzyga wyłącznie poprawność konserwatywną;
do jednoznacznego wyboru skoku potrzebna jest jeszcze
nierówność entropijna.
\end{proof}

\begin{przyklad}[Skok Burgersa]
Dla $f(u)=u^2/2$ i stałych stanów po obu stronach
fali $u_L\ne u_R$ warunek \eqref{eq:pde2-rh} daje
$\sigma=s'=(u_L+u_R)/2$.
Przy $u_L>u_R$ zachodzi $u_L>\sigma>u_R$:
charakterystyki z obu stron wpadają w falę uderzeniową.
Jeżeli $u_L<u_R$, taka nieciągłość spełnia wprawdzie
warunek Rankine'a--Hugoniota, ale nie jest dopuszczalnym
rozwiązaniem entropijnym.
\end{przyklad}

\begin{przyklad}[Fala rozrzedzenia]
\label{ex:pde2-rarefaction}
Dla równania Burgersa z danymi Riemanna
$u_0(x)=u_L$ dla $x<0$ oraz $u_0(x)=u_R$ dla $x>0$,
przy $u_L<u_R$, funkcja
\[
 u(t,x)=
 \begin{cases}
 u_L,&x\leq u_Lt,\\
 x/t,&u_Lt<x<u_Rt,\\
 u_R,&x\geq u_Rt
 \end{cases}
 \qquad(t>0)
\]
jest ciągła na każdym przekroju $t>0$ i spełnia
równanie Burgersa klasycznie w każdym z trzech
otwartych obszarów. Bezpośrednio sprawdzamy, że
w środkowej części $u_t=-x/t^2$, $u_x=1/t$,
zatem $u_t+uu_x=0$. Ciągłość na liniach
$x=u_Lt$ i $x=u_Rt$ sprawia, że oba skoki, wartości $u$
i strumienia $f(u)=u^2/2$, są zerowe.
Dla $\varphi\in C_c^\infty([0,T)\times\R)$ oraz
$0<\varepsilon<T$ stosujemy całkowanie przez części
na trzech częściach zbioru $(\varepsilon,T)\times\R$.
Wkłady na obu liniach wewnętrznych znoszą się,
a przy $t=T$ i na przestrzennych krańcach nośnika
test znika. Na dolnej krawędzi pozostaje
\[
 \int_\varepsilon^T\!\int_\R
 (u\varphi_t+f(u)\varphi_x)\,
 \mathrm d\lambda_1(x)\,\mathrm d\lambda_1(t)
 +\int_\R u(\varepsilon,x)\varphi(\varepsilon,x)\,
 \mathrm d\lambda_1(x)=0.
\]
Niech $M=\max\{|u_L|,|u_R|\}$. Zarówno $u(t,\cdot)$,
jak i $u_0$ są ograniczone w module przez $M$.
Poza przedziałem
$[\min(0,u_Lt),\max(0,u_Rt)]$ funkcje te są równe.
Stąd dla $t>0$
\[
 \int_\R |u(t,x)-u_0(x)|\,\mathrm d\lambda_1(x)
 \le 2Mt\bigl(\max(0,u_R)-\min(0,u_L)\bigr).
\]
W szczególności zachodzi zbieżność do danych początkowych
w $L^1_{\mathrm{loc}}$. W całce brzegowej możemy więc
zastąpić $u(\varepsilon,\cdot)$ przez $u_0$ z błędem
dążącym do zera, a $\varphi(\varepsilon,\cdot)$ dąży
jednostajnie do $\varphi(0,\cdot)$ na wspólnym zwartym
nośniku przestrzennym. W całce podwójnej korzystamy
z twierdzenia o zbieżności zdominowanej, ponieważ
$u$ oraz $f(u)$ są ograniczone, a test i jego pochodne
mają zwarty nośnik. Przejście z $\varepsilon$ do zera
daje dokładnie \eqref{eq:pde2-weak-form}.
To rozwiązanie jest entropijne; precyzyjną definicję
warunków entropijnych podamy w rozdziale o prawach zachowania.
\end{przyklad}

\begin{uwaga}[Dlaczego samo rozwiązanie słabe nie wystarcza?]
Dla danych z przykładu \ref{ex:pde2-rarefaction}
można też skonstruować skok poruszający się z prędkością
$(u_L+u_R)/2$. Spełnia on ten sam całkowy warunek
\eqref{eq:pde2-weak-form}, lecz narusza warunek entropijny.
Dla tego samego zagadnienia mamy więc co najmniej
dwa różne rozwiązania słabe. Selekcja rozwiązania
fizycznego wymaga dodatkowej informacji.
\end{uwaga}

\section{Równania w pełni nieliniowe i układ Charpita}
\label{sec:pde2-charpit}

Dla ogólnego równania pierwszego rzędu
\begin{equation}\label{eq:pde2-general-F}
 F(x,u(x),\nabla u(x))=0,
\end{equation}
gdzie $x\in\Omega\subset\R^n$,
$F\colon\Omega\times\R\times\R^n\to\R$,
rozpatrujemy dodatkowe zmienne $p\in\R^n$
reprezentujące gradient. Pojawia się
\emph{układ charakterystyczny Charpita}:
\begin{equation}\label{eq:pde2-charpit-system}
 \dot X=F_p(X,U,P),\qquad
 \dot U=P\cdot F_p(X,U,P),\qquad
 \dot P=-F_x(X,U,P)-F_u(X,U,P)P.
\end{equation}
Tutaj $F_x,F_p$ są gradientami po odpowiednich
zmiennych, a $F_u$ jest pochodną skalarną.

\begin{lemat}[Zachowanie równania i warunku styczności]
\label{lem:pde2-contact}
Niech $F\in C^2$ i niech $(X,U,P)$ będzie rozwiązaniem
\eqref{eq:pde2-charpit-system} na przedziale $I$.
Wtedy $s\mapsto F(X(s),U(s),P(s))$ jest stała.
Jeśli rozwiązania zależą klasy $C^1$ od parametru
$y\in V\subset\R^{n-1}$ i $F(X,U,P)=0$, to
\begin{equation}\label{eq:pde2-contact-ode}
 \frac{\mathrm d}{\mathrm ds}
 \bigl(U_{y_j}-P\cdot X_{y_j}\bigr)
 =-F_u(X,U,P)\bigl(U_{y_j}-P\cdot X_{y_j}\bigr)
\end{equation}
dla każdego $j=1,\ldots,n-1$.
\end{lemat}
\begin{proof}
Reguła łańcuchowa daje
\[
 \frac{\mathrm d}{\mathrm ds}F(X,U,P)
 =F_x\cdot F_p+F_u(P\cdot F_p)
 +F_p\cdot(-F_x-F_uP)=0.
\]
Różniczkując $R_j=U_{y_j}-P\cdot X_{y_j}$ po $s$,
wykorzystując \eqref{eq:pde2-charpit-system}, otrzymujemy
\[
 \dot R_j=P_{y_j}\cdot F_p+(F_x+F_uP)\cdot X_{y_j}.
\]
Ponieważ $F(X,U,P)=0$ dla każdego $(s,y)$,
mamy także
$F_x\cdot X_{y_j}+F_u U_{y_j}+F_p\cdot P_{y_j}=0$.
Po podstawieniu do wzoru na $\dot R_j$ dostajemy
$\dot R_j=-F_u(U_{y_j}-P\cdot X_{y_j})$, czyli
\eqref{eq:pde2-contact-ode}.
\end{proof}

\begin{twierdzenie}[Lokalna konstrukcja Charpita]
\label{thm:pde2-charpit}
Niech $n\geq2$, $\Omega\subset\R^n$ będzie otwarta,
$F\in C^2(\Omega\times\R\times\R^n)$,
$V\subset\R^{n-1}$ otwarty oraz
$\gamma\in C^2(V;\Omega)$ regularną parametryzacją
hiperpowierzchni. Niech $g\in C^2(V)$ i
$p_0\in C^1(V;\R^n)$ spełniają dla wszystkich $y\in V$
\begin{equation}\label{eq:pde2-charpit-initial}
 F(\gamma(y),g(y),p_0(y))=0,\qquad
 p_0(y)\cdot\partial_{y_j}\gamma(y)=\partial_{y_j}g(y)
 \quad(j=1,\ldots,n-1).
\end{equation}
Jeżeli w $y_0\in V$ wektor
$F_p(\gamma(y_0),g(y_0),p_0(y_0))$
nie należy do $T_{\gamma(y_0)}\Sigma$,
gdzie $\Sigma=\gamma(V)$ lokalnie, to istnieje
otoczenie $U$ punktu $\gamma(y_0)$ i rozwiązanie
$u\in C^2(U)$ równania \eqref{eq:pde2-general-F}
przyjmujące dane $u(\gamma(y))=g(y)$ w pobliżu $y_0$.
Jego gradient na powierzchni danych jest równy $p_0$.
W klasie rozwiązań $C^2$, które realizują te same
dane $g,p_0$, konstrukcja jest lokalnie jednoznaczna.
\end{twierdzenie}
\begin{proof}
Rozwiążmy układ \eqref{eq:pde2-charpit-system}
z danymi $(X,U,P)(0,y)=(\gamma(y),g(y),p_0(y))$.
Zależność rozwiązań ODE od parametrów daje funkcje
klasy $C^1$ względem $(s,y)$ w pobliżu $(0,y_0)$.
Lemat \ref{lem:pde2-contact} i pierwszy warunek
\eqref{eq:pde2-charpit-initial} pokazują, że
$F(X,U,P)=0$; drugi warunek i równanie
\eqref{eq:pde2-contact-ode} dają
$U_{y_j}=P\cdot X_{y_j}$ przez jednoznaczność
skalarnego równania liniowego. Ponadto z samego
układu wynika $U_s=P\cdot X_s$.

W punkcie $(0,y_0)$ kolumny macierzy $D_{(s,y)}X$
są równe $F_p$ i $\partial_{y_j}\gamma$;
warunek poprzeczności zapewnia, że macierz ta jest
odwracalna. Z twierdzenia o funkcji odwrotnej $X$
stanowi lokalny układ współrzędnych w $\Omega$.
Definiujemy $u(X(s,y))=U(s,y)$. Wspomniane wcześniej
równości są równoważne tożsamości
$\mathrm dU=P\cdot\mathrm dX$, zatem
$\nabla u(X(s,y))=P(s,y)$. Ponieważ $P$ i lokalna
odwrotność $X^{-1}$ są klasy $C^1$, gradient $\nabla u$
jest klasy $C^1$ na $U$, a więc $u\in C^2(U)$.
Ponieważ $F(X,U,P)=0$, otrzymujemy równanie
\eqref{eq:pde2-general-F}.

Jeśli $u\in C^2$ jest innym rozwiązaniem z tym samym
gradientem początkowym, to wzdłuż odpowiednich
charakterystyk jego wartości i pochodne spełniają
\eqref{eq:pde2-charpit-system}: wystarczy zróżniczkować
$F(x,u,\nabla u)=0$ i skorzystać z symetrii hesjanu.
Jednoznaczność ODE implikuje zgodność obu rozwiązań
na wspólnym otoczeniu punktu początkowego.
\end{proof}

\begin{przyklad}[Eikonal i wybór gałęzi]
Niech $F(x,u,p)=|p|^2-1$ w $\R^2$.
Na prostej $\Sigma=\{(x,0):x\in\R\}$ zadamy
$u(x,0)=0$. Zgodność gradientu stycznego wymusza
$p_x=0$, a równanie wymusza $p_y=\pm1$.
Są zatem dwie dopuszczalne gałęzie danych gradientu,
które prowadzą do rozwiązań $u(x,y)=y$ i $u(x,y)=-y$.
Oba są niecharakterystyczne dla odpowiedniego wektora
$F_p=2p$. Samo zadanie wartości na prostej nie określa
więc jednoznacznie rozwiązania w pełni nieliniowego.
\end{przyklad}

\section{Równanie Hamiltona--Jacobiego i charakterystyki Hamiltona}
\label{sec:pde2-hj}

Ważną szczególną postacią równania pierwszego rzędu jest
\begin{equation}\label{eq:pde2-hj}
 u_t+H(x,\nabla_xu)=0,
 \qquad u(0,x)=u_0(x).
\end{equation}
W tej sekcji $H\in C^2(\R^d\times\R^d)$ nie zależy
jawnie od czasu ani od wartości $u$.

\begin{stwierdzenie}[Metoda charakterystyk Hamiltona]
\label{prop:pde2-hamilton}
Niech $u_0\in C^2(\R^d)$, $H\in C^2(\R^d\times\R^d)$
i niech $\xi_0\in\R^d$. Rozwiążmy lokalnie układ
\begin{equation}\label{eq:pde2-hamilton-odes}
 \dot X=H_p(X,P),\qquad
 \dot P=-H_x(X,P),\qquad
 \dot U=P\cdot H_p(X,P)-H(X,P),
\end{equation}
z warunkami $X(0;\xi)=\xi$,
$P(0;\xi)=\nabla u_0(\xi)$ i $U(0;\xi)=u_0(\xi)$.
Dla dostatecznie małych $|t|$ i $\xi$ bliskich $\xi_0$
mapa $\xi\mapsto X(t;\xi)$ jest lokalnie odwracalna.
W tych współrzędnych
$u(t,X(t;\xi))=U(t;\xi)$ definiuje rozwiązanie
klasy $C^1$ równania \eqref{eq:pde2-hj}, a ponadto
$\nabla_xu(t,X(t;\xi))=P(t;\xi)$.
\end{stwierdzenie}
\begin{proof}
Istnienie ODE i lokalna odwracalność wynikają jak
w twierdzeniu \ref{thm:pde2-quasilinear}, ponieważ
$D_\xi X(0;\xi)=I$. Ustalmy współrzędną $\xi_j$
i oznaczmy $R_j=U_{\xi_j}-P\cdot X_{\xi_j}$.
Korzystając z \eqref{eq:pde2-hamilton-odes} oraz
reguły łańcuchowej, obliczamy
\begin{align*}
 \dot R_j&=\partial_{\xi_j}(P\cdot H_p-H)
 -\dot P\cdot X_{\xi_j}-P\cdot\partial_{\xi_j}H_p\\
 &=P_{\xi_j}\cdot H_p-H_x\cdot X_{\xi_j}
 -H_p\cdot P_{\xi_j}+H_x\cdot X_{\xi_j}=0.
\end{align*}
Początkowo
$R_j(0,\xi)=\partial_{\xi_j}u_0-\nabla u_0\cdot e_j=0$,
zatem $R_j=0$. Wynika stąd
$\nabla_xu(t,X(t;\xi))=P(t;\xi)$. Z kolei pochodna
$u(t,X(t;\xi))=U(t;\xi)$ po $t$ daje
$u_t+\nabla_xu\cdot H_p=P\cdot H_p-H$, a więc
$u_t+H=0$, jak należało wykazać.
\end{proof}

\begin{przyklad}[Hamiltonian kinetyczny i powstawanie kaustyki]
\label{ex:pde2-hj-blowup}
Niech $d=1$ i $H(p)=p^2/2$. Wtedy
$\dot P=0$, $\dot X=P$ oraz $\dot U=P^2/2$.
W konsekwencji
\[
 X(t;\xi)=\xi+t u_0'(\xi),\qquad
 P(t;\xi)=u_0'(\xi),\qquad
 U(t;\xi)=u_0(\xi)+\frac t2|u_0'(\xi)|^2.
\]
Dla $u_0(x)=-x^2/2$ otrzymujemy
$X(t;\xi)=(1-t)\xi$. Jeśli $0\leq t<1$, to
\[
 u(t,x)=-\frac{x^2}{2(1-t)},\qquad
 u_x(t,x)=-\frac{x}{1-t}.
\]
W chwili $t=1$ mapa $\xi\mapsto X(t;\xi)$ traci
odwracalność, a dla ustalonego $x\ne0$ gradient
rozwiązania staje się nieograniczony, gdy $t\uparrow1$.
Zjawisko utraty odwracalności rzutu
$(X,P)\mapsto X$ nazywa się w tym kontekście
powstawaniem \emph{kaustyki}. Rozwiązania lepkościowe
umożliwią później kontynuację odpowiednich zadań
Hamiltona--Jacobiego poza pierwszą kaustyką.
\end{przyklad}

\section{Podsumowanie: kiedy metoda charakterystyk działa?}
\label{sec:pde2-summary}

Metoda charakterystyk wykorzystuje zależność
\eqref{eq:pde2-chain}, ale jej wynik zależy od struktury
równania:
\begin{enumerate}
\item Dla transportu liniowego z wystarczająco regularnym
polem prędkości charakterystyki tworzą przepływ
dyfeomorfizmów i zapewniają globalne rozwiązanie klasyczne
na skończonym przedziale czasu.
\item Dla równań quasiliniowych lokalne rozwiązanie istnieje,
dopóki układ charakterystyczny daje odwracalną mapę
początkowych etykiet $\xi$ na położenia $x$.
\item Dla praw zachowania po utracie gładkości naturalne
stają się rozwiązania słabe. Sama słaba postać równania
nie wystarcza zwykle do jednoznaczności; potrzebne są
warunki entropijne.
\item Dla równań w pełni nieliniowych trzeba śledzić
również gradient $p=\nabla u$. Warunek poprzeczności
$F_p$ względem powierzchni danych jest kluczowy,
a rozwiązanie może istnieć jedynie lokalnie nawet
przy bardzo gładkich danych.
\end{enumerate}

W rozdziale następnym przejdziemy do równania Laplace'a,
teorii funkcji harmonicznych i zagadnień Dirichleta.

